\section{问题重述}

\subsection{问题背景}

低温作业人员通过三层防护服抵御严寒，其中功能层含相变材料，其放热特性由附件DSC数据描述。题目要求在给定衣料热物性与人体参数的基础上，建立人体—防护服—环境瞬态传热模型，并进一步完成隔热层增厚与功能材料强化设计，整体背景如图~\ref{fig:problem-background}所示。

\begin{figure}[!htbp]
  \centering
  \includegraphics[width=\FigureMechanismWidth]{\ProjectRoot/modules/10_restatement/paper/问题背景.png}
  \caption{问题背景示意图}
  \label{fig:problem-background}
\end{figure}
\FloatBarrier

\subsection{问题提出}

针对低温防护服的传热预测与优化设计问题，具体研究的问题如下：

\indent\textbf{问题一：}在$-40\,^{\circ}\mathrm C$、静止无风条件下建立三层防护服瞬态传热模型，预测贴身侧温度首次降至$15\,^{\circ}\mathrm C$和$10\,^{\circ}\mathrm C$的时刻。

\indent\textbf{问题二：}将外部风速改为$3\,\mathrm{m/s}$、人体状态改为轻微运动，重新确定内外换热边界并计算两级温度阈值对应的返回时间。

\indent\textbf{问题三：}在材料费用与随站立时间衰减的承重约束下增加最外隔热层厚度，以最大化实验者的有效持续时间并确定最优增厚方案。

\indent\textbf{问题四：}以问题三最优持续时间为目标，在原始厚度和材料支出不变的条件下，反求功能层单位质量放热能力所需的最小提高比例。

\FloatBarrier
